\section{Introduction}\label{sec:introduction}

For a finite set \(P\subset\mathbb R^d\), write
\[
 \Delta(P)=\{|p-q|:p,q\in P,\ p\ne q\}.
\]
The distinct-distances problem asks how small \(|\Delta(P)|\) can be
among all sets of a prescribed cardinality. Erd\H{o}s introduced the
problem in 1946~\cite{erdos1946}. It relates a simple metric statistic
to the geometric structure of a finite set: when distances are few,
many pairs repeat their lengths, but the arrangement of those repetitions is
constrained by Euclidean geometry.
The higher-dimensional conjecture predicts a constant multiple of
\(n^{2/d}\) distances among any \(n\) points of \(\mathbb R^d\), for
every fixed \(d\ge3\); see, for example,
\cite{BardwellEvansSheffer2019,tidoryuzakharov2026}.
We prove this prediction.

\begin{theorem}\label{thm:main}
For every integer \(d\ge3\), there is a constant \(c_d>0\) such that
every finite set \(P\subset\mathbb R^d\) of \(n\ge2\) distinct points
satisfies
\[
 |\Delta(P)|\ge c_d n^{2/d}.
\]
\end{theorem}

This resolves the higher-dimensional Erd\H{o}s distinct-distances
conjecture positively. No assumption is made on the position, spacing,
or concentration of the points. The power of \(n\) is optimal:
the integer grid \(\{1,\ldots,t\}^d\) has \(t^d\) points, and its squared
distances are integers between \(1\) and \(d(t-1)^2\). In the plane,
the square-grid construction gives \(O(n/\sqrt{\log n})\)
distances~\cite{erdos1946}; the planar problem consequently
has a different conjectural order.

\subsection{Historical context and the methods of the proof}

Higher-dimensional distance estimates developed in part through
incidence bounds for curves and spheres. Clarkson, Edelsbrunner,
Guibas, Sharir, and Welzl bounded the number of pairs realizing a fixed
positive distance in three dimensions, giving a distinct-distance
lower bound at the square-root scale, up to a slowly growing
factor~\cite{ClarksonEdelsbrunnerGuibasSharirWelzl1990}.
Aronov, Pach, Sharir, and Tardos improved the three-dimensional lower
bound to \(\Omega_\varepsilon(n^{77/141-\varepsilon})\) for each
\(\varepsilon>0\) and propagated improvements to higher
dimensions~\cite{AronovPachSharirTardos2004}.
Their three-dimensional estimate is \emph{pinned}: it finds a point
of the given set from which that many distinct distances occur.

Solymosi and Vu developed two recursive inequalities connecting
distance bounds in lower and higher dimensions
\cite[Theorems~2.1--2.2]{SolymosiVu2008}.
The recurrences have separate forms for the total number of distances
and for the maximum number measured from one point of the set. We use
the forms for total distances in the comparisons below. Their original theorem gives
\(\Omega(n^{0.5643})\) distances in dimension three and
\(\Omega_d(n^{2/d-2/(d(d+2))})\) for every fixed \(d\ge4\)
\cite[Theorem~1.1]{SolymosiVu2008}.

A different approach counts repeated distances through rigid motions
matching subsets of the point sets. Elekes and Sharir developed this viewpoint for the planar
problem~\cite{elekessharir2011}. Guth and Katz reduced the resulting
planar motion count to incidences of points and lines in three-dimensional
space and combined polynomial partitioning with ruled-surface geometry to prove
the planar bound \(\Omega(n/\log n)\)
\cite[Theorem~1.1 and Sections~2--4]{guthkatz2015}.
That bound supplies the two-dimensional input to our induction.
Bardwell-Evans and Sheffer extended the rigid-motion approach to
higher dimensions through a reduction to incidences of
\((d-1)\)-flats in \(\mathbb R^{2d-1}\)
\cite{BardwellEvansSheffer2019}. Their work identifies structural
restrictions on the resulting flats and the need to control their
concentration in algebraic varieties.

Combining the Guth--Katz planar bound with the first Solymosi--Vu
recurrence gives \(\Omega(n^{3/5}/(\log n)^{2/5})\) distances in
dimension three~\cite[Theorem~2.1]{SolymosiVu2008}.
More recently, Tidor, Yu, and Zakharov proved
\(n^{2/3-o(1)}\) distinct distances for arbitrary point sets in
\(\mathbb R^3\)~\cite[Theorem~1.1]{tidoryuzakharov2026}.
Inserting their three-dimensional distance bound into the Solymosi--Vu
recurrences gives \(n^{8/17-o(1)}\) in dimension four and
\(n^{3/8-o(1)}\) in dimension five
\cite[Corollaries~2.3 and~2.6]{SolymosiVu2008}.
The \(o(1)\) terms here absorb the subpolynomial loss in that input.
The leading exponents in these two deductions still fall below the
conjectured value \(2/d\).
Bounds for points on prescribed surfaces, or for distances from a
single point, have different hypotheses or conclusions.\footnote{An
earlier manuscript of Aksoy Yazici claimed
the constant-factor bound in all dimensions \(d\ge3\), but was
withdrawn by the author because the proof of its Theorem~1.2 is not
correct~\cite{AksoyYazici2020Withdrawn}.}

Tidor, Yu, and Zakharov's proof develops rigid-motion flat geometry,
regular approximate complete intersections, and degree-sensitive
incidence estimates. These methods provide close antecedents for
several parts of the argument below: the physical and
opposite-orientation flat families, the recovery of actual families
of contained flats from tangent data, and the geometric treatment of
projected curves
\cite[Sections~3--6 and~10.2--10.3]{tidoryuzakharov2026}.

The algebraic argument draws on Hilbert function estimates of Chardin
and Chardin--Philippon
\cite{chardin1989,chardinphilippon1999,chardinphilippon2002erratum},
and on Walsh's work on approximate complete intersections and
algebraic concentration~\cite{walsh2020,walsh2023}.
Hilbert functions measure how many independent
polynomial conditions a set can impose. Approximate complete
intersections cut out an irreducible variety as a component of the
common zero set of equations whose product of degrees is controlled
by the variety's degree.
We use these two ideas to keep track of concentration at several
polynomial degree scales. The uniform concentration theorem for
evaluation flats of skew forms, including its exceptional parameter
families, is proved in full in Appendix~\ref{conc:section}.
The real incidence arguments, including the planar distinct-distance
bound, are supplied in Appendix~\ref{ci:section}; their relation to
the work of Szemer\'edi--Trotter, Beck, Elekes--Sharir, and Guth--Katz
is recorded where each argument is used.

\subsection{Why the constant factor matters}

An estimate \(|\Delta(P)|\ge n^{2/d-o(1)}\) does not by itself imply
Theorem~\ref{thm:main}. The ratio \(|\Delta(P)|/n^{2/d}\) could still
tend to zero more slowly than every negative power of \(n\).
To obtain a constant factor, the proof must exclude every such
sequence, without first fixing a power saving.

We therefore argue in the least dimension \(d\ge3\) in which the
theorem could fail, and use the notation
\[
 N=|P|,\qquad B=N^{1/d},\qquad A=B^2,\qquad
 M=1+|\Delta(P)|.
\]
Failure would give a sequence with \(N\to\infty\) and \(M/A\to0\).
The lower-dimensional bounds limit the number of points in any
proper affine flat. They do not prevent substantial concentration
on more general varieties, nor do they identify a single polynomial
degree at which all of that concentration can be studied.
The multiscale argument below addresses this second difficulty.

\subsection{Overview of the proof}

\paragraph{1. Select directions that admit interpolation.}
For distinct points \(p,q\), let
\(\pi_p(q)=[q-p]\in\mathbb P^{d-1}\) denote their projective direction.
We seek a positive fraction of the unordered pairs of points such
that, at each center \(p\), homogeneous polynomials of degree
\(\lfloor A\rfloor\) can prescribe arbitrary values at the selected
displacement vectors \(q-p\). Equivalently, the corresponding evaluation
vectors are linearly independent. The selection must be symmetric:
a pair is retained or removed at both endpoints.

The principal geometric obstruction would be a collection \(C_p\)
of projective curves at each center, each curve spanning at most a
projective plane, with total degree \(o(N/A)\), such that almost
every pair \(\{p,q\}\) satisfies
\[
 \pi_p(q)\in C_p\quad\hbox{or}\quad\pi_q(p)\in C_q.
\]
The union of the corresponding lines through \(p\) is a cone;
the small total degree explains the term \emph{sparse cones}.
Theorem~\ref{sc:theorem} rules out this coverage, including the
choice of either endpoint for each pair. Once it is excluded, greedy
removal of curves containing too many directions leaves a positive
fraction of pairs with controlled concentration on every projective
variety. A maximal symmetric selection then gives the interpolation
property (Propositions~\ref{red:directions} and
\ref{red:independent-stars}). This last linear-algebra argument costs
only a fixed fraction of pairs.

\paragraph{2. Exclude sparse cones at all degree scales.}
The sparse-cones theorem occupies the central part of the paper.
Section~\ref{sec:profiles} divides the point set according to how
successive polynomial equations cut down the dimension of the
varieties containing its points. It tests degrees in windows
\([B e^{-s},B e^s]\) with several width parameters \(s\), instead of
fixing a single power of \(N\). Points isolated by equations of
sufficiently small degree form a core that is handled directly.
Each remaining piece carries a fixed chain of proper hypersurface
cuts, ending in isolated points.
A cut is proper when its equation does not vanish identically on
the component being cut. The testing degrees at these cuts form a
\emph{scale profile}. Between consecutive cuts, the dimension stays
constant, and Hilbert-function estimates give a lower bound on the
growth of polynomial restriction rank. B\'ezout's theorem bounds
the degrees and point counts of the intervening components.

The comparison uses labels \((p,\pi_p(q))\), which record both the
center and the direction of an ordered pair. At each fixed center,
sparse-cone coverage puts the directions on curves of small total
degree, limiting their polynomial rank. The profiles give an incompatible rank lower bound
for a positive fraction of pairs. A complication is that radial
projection can identify several components, or several points on
one component. Section~\ref{geo:section} controls this behavior by
secant geometry. For separated windows, Subsection~\ref{dep:section}
bounds the resulting loss of degree at the actual centers.

Two limiting cases leave no strict margin in the general rank
comparison. At the outermost degree window, common equations with
independent differentials on most components allow a direct geometric
comparison of their projected images. At the last one-dimensional
stage for narrower center windows, the argument identifies the projection
of the assigned terminal curve with a curve in the covering family.
A global estimate has already removed all such incidences at a
cost of \(o(N^2)\) pairs (Proposition~\ref{geo:terminal}). These two
arguments in Section~\ref{sec:sparse-cones} supply the endpoint
control needed for an arbitrarily slow limit \(M/A\to0\).

\paragraph{3. Turn equal distances into intersections of flats.}
Take a generic rotated copy \(Q\) of \(P\) and transport the selected
pairs to \(Q\). For \((p,q)\in P\times Q\),
consider the affine flat
\[
 F_{p,q}=\{(Z,c):Z(p+q)/2+c=(p-q)/2\},
 \qquad Z^{\mathsf T}=-Z,\quad c\in\mathbb C^d.
\]
The coordinates \((Z,c)\) identify with skew forms on
\(\mathbb C^{d+1}\) by adjoining \(c\) as the last column and
\(-c^{\mathsf T}\) as the last row, with last diagonal entry zero.
Join \((p,q)\) to \((p',q')\)
when both physical pairs have been selected in Step 1 and
\(|p-p'|=|q-q'|\). Cauchy--Schwarz gives
\(\Omega(N^4/M)\) edges. On \(F_{p,q}\), the intersection with a
neighbor is the flat defined by
\[
 Z(u+w)=u-w,\qquad u=p'-p,\quad w=q'-q.
\]
The interpolation property ensures that these local flats are
distinct and that their parameters have polynomial separators of
degree \(O(A)\): one parameter can be distinguished from all the
others by a polynomial of that degree.

Appendix~\ref{conc:section} proves concentration bounds for these
evaluation flats, both in the full skew-form space and inside one
fixed pair flat. In the local application an exceptional parameter
family corresponds to many matches under one rigid motion. For
\(d\ge4\), its required size bound follows from the total number of
points. For \(d=3\), Section~\ref{mot:section} bounds
\(\sum_g |gP\cap Q|^2\) over motions with at least \(CA\) matches
by \(O(N^4/A)\), with both orientations included. Removing the
corresponding edges costs \(o(N^4/M)\). After this removal and
pruning vertices of small degree, the minimum number of neighbors,
divided by \(A^{d-1}\), tends to infinity.
The local concentration bound says that
a nonzero polynomial of degree at most \(A\) on a pair flat can
vanish on at most \(O(A^{d-1})\) of its neighbor flats.

\paragraph{4. Sample once and compare polynomial ranks.}
The remaining contradiction is short. Sample the retained pair flats
independently with a small fixed probability. With positive probability,
the sample is small enough that the global concentration bound forces
the whole family to impose more independent polynomial conditions
in degree \(\lfloor A\rfloor\) than the sample. Simultaneously,
every retained flat has more than \(CA^{d-1}\) sampled neighbors.
Thus a polynomial of degree at most \(\lfloor A\rfloor\) vanishes on every sampled flat
but is nonzero on at least one retained flat \(F\).
The restriction of the polynomial to
\(F\) then violates local concentration. The sample uses one fixed
probability, and all concentration constants are uniform; this is
where the preceding quantitative estimates produce the positive
constant in Theorem~\ref{thm:main}.

\paragraph{Organization and prerequisites.}
Section~\ref{sec:reduction} carries out Steps 1, 3, and 4, assuming
the sparse-cones, concentration, and rich-motion theorems. It can
therefore be read first as the formal assembly of the argument.
Section~\ref{poly:section} proves the polynomial interpolation,
regularity, and multiplicity estimates. Sections~\ref{geo:section}
through~\ref{sec:sparse-cones} establish the geometric and multiscale
arguments of Step 2. Section~\ref{mot:section} proves the rich-motion
estimate, and Appendices~\ref{conc:section} and~\ref{ci:section}
provide the concentration and real incidence foundations.
Figure~\ref{fig:proof-map} summarizes the principal dependencies.
We use characteristic-zero algebraic geometry, including projective
degree and B\'ezout's theorem, generic linear sections, regular local
rings and Koszul resolutions, and projective-bundle intersection
formulas. The precise forms needed are specified where they are
used. Constants may depend on the fixed ambient dimension; any
additional dependence is indicated explicitly.

\begin{figure}[htbp]
\centering
\begin{tikzpicture}[
 box/.style={draw=linkblue!65,rounded corners=2pt,fill=figuregray,
             align=center,font=\small,text width=4.65cm,
             minimum height=1.15cm,inner sep=5pt},
 application/.style={align=center,font=\small,text width=4.1cm,
                     inner sep=3pt},
 arrow/.style={-{Latex[length=2mm]},thick,draw=linkblue!80}]
 \node[box] (poly) at (-3.2,0.65)
   {Polynomial estimates\\and scale profiles\\
    Sections~\ref{poly:section} and~\ref{sec:profiles}};
 \node[box] (geo) at (3.2,0.65)
   {Secant geometry\\and terminal-curve deletion\\
    Section~\ref{geo:section}};
 \node[box,text width=8.8cm] (sparse) at (0,-1.6)
   {\textbf{Sparse-cones argument} (Section~\ref{sec:sparse-cones})\\
    Common windows: rank and outer endpoint\\
    Separated windows: rank and final curve endpoint,\\
    using projection depth at actual centers};
 \node[box] (stars) at (0,-3.8)
   {Symmetric direction selection\\with independent evaluations\\
    Proposition~\ref{red:independent-stars}};
 \node[box] (graph) at (0,-5.7)
   {Equal-distance graph\\and its pair flats\\
    Section~\ref{sec:reduction}};
 \node[box,text width=3.8cm] (conc) at (-5.3,-5.7)
   {Uniform concentration\\of evaluation flats\\
    Appendix~\ref{conc:section}};
 \node[box,text width=3.25cm] (motions) at (5.4,-5.7)
   {Very rich motions\\in dimension three\\
    Section~\ref{mot:section}};

 \node[draw=linkblue!65,rounded corners=2pt,fill=figuregray,
       minimum width=9cm,minimum height=2.4cm,inner sep=0pt]
       (applications) at (0,-9) {};
 \node[font=\small\bfseries,align=center] at (0,-8.08)
   {Applications to the retained pair flats};
 \node[application] (global) at (-2.25,-9.2)
   {Global concentration\\and the Hilbert bound:\\rank lower bound\\
    Lemmas~\ref{red:global-concentration} and~\ref{poly:hilbert}};
 \node[application] (local) at (2.25,-9.2)
   {Local concentration:\\cap on neighbor flats in a\\
    proper low-degree zero set\\
    Lemma~\ref{red:local-concentration}};
 \draw[linkblue!25] (0,-8.55) -- (0,-9.94);

 \node[box,text width=5.8cm,fill=figureblue] (end) at (0,-11.45)
   {One random sample\\with fixed inclusion probability\\
    and a rank contradiction\\
    Proposition~\ref{red:main-contradiction}};

 \draw[arrow] (poly.south) -- ([xshift=-3.2cm]sparse.north);
 \draw[arrow] (geo.south) -- ([xshift=3.2cm]sparse.north);
 \draw[arrow] (sparse.south) -- (stars.north);
 \draw[arrow] (stars.south) -- (graph.north);
 \draw[arrow] (graph.south) -- (applications.north);
 \draw[arrow] (conc.south) -- (applications.north west);
 \draw[arrow] (motions.south) --
   node[pos=.48,right,xshift=3pt,align=left,font=\footnotesize,inner sep=3pt]
        {\(d=3\) deletion\\for the local\\isotropic cap}
   (local.north east);
 \draw[arrow] (applications.south) -- (end.north);
\end{tikzpicture}

\caption{Principal dependencies. Both endpoint arguments belong to the
sparse-cones theorem. The separated-window argument uses projection depth
at actual centers. Global concentration combines with the component
Hilbert bound; the rich-motion bound permits the \(d=3\) edge deletion
that supplies the local isotropic cap. The contradiction uses one sample with a
fixed inclusion probability. Appendix~\ref{ci:section} supplies the planar
induction base and the real incidence inputs to the geometric and motion
arguments.}
\label{fig:proof-map}
\end{figure}
